\chapter*{Introduction to Part~III}
\addcontentsline{toc}{chapter}{Introduction to Part~III}
\label{ch:part3-intro}

In this Part, I turn to neutrinos and free-streaming radiation --- species that, like dark matter, interact too feebly with ordinary matter to be easily probed in terrestrial experiments, yet leave distinct and measurable imprints on cosmic evolution. A particularly clean and robust signature is the neutrino-induced phase shift in the acoustic oscillations of the primordial plasma~\cite{Bashinsky:2003tk}. After decoupling from the photon-baryon fluid, roughly a second after the Big Bang at temperatures of about 1 MeV, neutrinos propagate freely at nearly the speed of light, outrunning the sound waves in the primordial plasma. Their gravitational influence pulls these acoustic oscillations slightly ahead of the sound horizon, producing a characteristic and coherent phase shift that remains imprinted in cosmological observables --- most directly in the anisotropies of the CMB, but also in large-scale structure tracers that inherit the same acoustic patterns.  Most importantly, this shift carries a specific scale dependence that is uniquely difficult to mimic through modifications to initial conditions or other cosmological parameters, making it a pristine observable to look for in order to probe the free-streaming nature of neutrinos and, more broadly, of any light relic species in the early Universe~\cite{Baumann:2015rya}.

Isolating this signature is especially important given that any suppression or modification of the phase shift would signal non-standard neutrino interactions --- either among themselves or with a dark sector --- which are ubiquitous in many extensions of the Standard Model and have been shown to potentially alleviate persistent cosmological issues, such as the Hubble tension~(e.g.~\cite{Ghosh:2019tab}). Building on the substantial body of theoretical and observational work developed over the past decade (see e.g.~\cite{Wallisch:2018rzj, Montefalcone:2025ibh} and references therein) the three chapters in this Part develop a comprehensive analysis framework for robustly isolating this free-streaming signature across multiple cosmological probes.

In Chapter~\ref{ch:neutrino-phase-shift}, I introduce two complementary template-based methods for extracting the phase shift from CMB power spectra, and apply them to current data from \textit{Planck}~\cite{Planck:2018vyg}, the Atacama Cosmology Telescope (ACT)~\cite{ACT:2025fju}, and the South Pole Telescope~\cite{SPT-3G:2025bzu}, obtaining a high-significance detection consistent with the Standard Model prediction of three free-streaming neutrino species.

In Chapter~\ref{ch:neutrino-interactions}, I extend this framework to a well-motivated class of interacting neutrino scenarios and derive a direct mapping between the measured phase-shift amplitude and the epoch at which neutrinos began free-streaming. The resulting constraints on non-standard neutrino interactions are competitive with existing bounds, and are obtained from a single, robust observable without requiring dedicated model-dependent analyses.


Finally, in Chapter~\ref{ch:neutrino-21cm}, I explore for the first time how this signature manifests in the 21-cm power spectrum from cosmic dawn, extending the reach of existing CMB and BAO measurements from galaxy surveys. The 21-cm signal carries a qualitatively richer acoustic structure than these earlier probes, owing to the presence of velocity-induced acoustic oscillations (VAOs) --- additional oscillatory features sourced by the supersonic relative motion between baryons and dark matter --- which imprint a phase response that is distinct from that of the standard baryon acoustic oscillations (BAOs). By characterizing both contributions, I demonstrate that the neutrino-induced phase shift leaves a distinct and trackable imprint in the 21-cm spectrum, laying the foundation for its extraction from ongoing and upcoming observational campaigns.
